% Countably many independent coefficient directions for the descent module.
\subsection{Prime-indexed harmonic defects and infinitely many descent generators}

Use the signed braid representation of R1, E.65.8 and E.65.15.
For each odd prime $p$, let
\[
 \vartheta_p:F\longrightarrow\operatorname{SL}_2(\mathbb F_p),
 \qquad x\longmapsto
 \begin{pmatrix}1&2\\0&1\end{pmatrix},\quad
 y\longmapsto\begin{pmatrix}1&0\\-2&1\end{pmatrix}.
 \tag{B.AC1}
\]
For $h\in H$, let $\mathcal V_{p,h}$ be the reduction of
the unique upper-unipotent intertwiner in R1 E.65.15, and put
\[
 W_{p,h}=\vartheta_p(\beta_h)^{-1}\mathcal V_{p,h},\qquad
 T_p=C_{\operatorname{SL}_2(\mathbb F_p)}(S).
 \tag{B.AC2}
\]
Write $n_p=p-\left(\frac{-1}{p}\right)$.  The group $T_p$ is
cyclic of order $n_p$, and its quadratic character is specified by
\[
 \varepsilon_p:T_p\longrightarrow\mathbb F_2,\qquad
             M^{n_p/2}=(-I)^{\varepsilon_p(M)}.
 \tag{B.AC3}
\]
Define on the actual multiplier-one group
\[
 w_p(h)=\varepsilon_p(W_{p,h}),\qquad
                   d_p(h)=w_p(h)+w_p(Dh).
 \tag{B.AC4}
\]

\begin{result}{B.2.30}{Proved}{Countably many independent coefficient directions and their full descent shifts}
For every odd prime $p$, (B.AC4) gives continuous
characters of $H$, invariant under conjugation by $G$.
There is a unique full continuous character
$\overline d_p:G\to\mathbb F_2$ restricting to $d_p$
and vanishing on $G_{\mathbb Q}$.  It vanishes on
full $\widehat{GT}$ and on full $C_L$.

Let $p_1<p_2<\cdots$ be the primes at least five which
are congruent to three or five modulo eight, and set
\[
 \kappa_0=\Delta,\qquad
 \kappa_1=\overline d_2,\qquad
               \kappa_{i+1}=\overline d_{p_i}\quad(i\geq1).
 \tag{B.AC5}
\]
Here $\overline d_2$ retains its earlier harmonic-defect
notation from B.2.26; the new subscripts $p_i$ are odd primes.
The restrictions of all $\kappa_j$ to $\ker D$ are
linearly independent.  With $t$ acting by $D^*$ on $V$,
the literal map
\[
 \bigoplus_{j\geq0}\mathbb F_2[t]\longrightarrow V,
               \qquad(P_j)_j\longmapsto
                         \sum_j P_j(D^*)\kappa_j
 \tag{B.AC6}
\]
is injective.  Moreover
\[
 \dim_{\mathbb F_2}(V/D^*V)=\aleph_0,
 \tag{B.AC7}
\]
so no finite family of characters and all of its descent
iterates can generate $V$.

The specified continuous homomorphism
\[
 \Xi_\infty:G\longrightarrow
             \prod_{(j,n)\in\mathbb N_0^2}\mathbb F_2,\qquad
 g\longmapsto\bigl(\kappa_j(D^ng)\bigr)_{j,n}
 \tag{B.AC8}
\]
is onto already on $H$ and vanishes on $\widehat{GT}$.
Its transition under $D$ shifts the second index.
Every finite set of coordinates has full gentle image
and trivial GT image.  For every $j\geq1$ and $n\geq0$,
the coordinate $\kappa_jD^n$ also vanishes on full $C_L$.
\end{result}
\begin{proof}
We first verify all source normalizations in (B.AC2).
Let $\widetilde s=uvu$ and $z=(uv)^3$ in the actual
central braid extension used in R1.  The multiplier-one
automorphism fixes $z$ and preserves braid abelianization.
Its lift of the source relation
$\Phi_h(s)=\beta_hs\beta_h^{-1}$ can differ from
$\beta_h\widetilde s\beta_h^{-1}$ only by $z^a$.
Abelianization gives $3=3+6a$ in $\widehat{\mathbb Z}$,
hence $a=0$.  Thus the signed matrix relation holds,
and the covariance of R1 E.65.15 gives
$W_{p,h}S=SW_{p,h}$ exactly.  The possible simultaneous
choice of $S$ or $-S$ has the same centralizer.
The identical argument for the cubic braid lift uses
its abelian value two and again gives $6a=0$;
thus the pure cubic-frame identity is signed as well.

Solving this matrix commutation equation gives
\[
 T_p=\{aI+bS:a^2+b^2=1\}.
 \tag{B.AC9}
\]
When $-1$ is a square, the two eigenvalues identify
this group with $\mathbb F_p^\times$.  Otherwise it
is the norm-one subgroup of $\mathbb F_{p^2}^\times$.
This proves its cyclicity and stated order.
In both cases $4$ divides $n_p$, so
$\varepsilon_p(-I)=0$ and
\[
 \varepsilon_p(S)=n_p/4\bmod2
 =1\quad\text{for }p\equiv3,5\pmod8.
 \tag{B.AC10}
\]
These elementary matrix assertions and the source
matrix (B.AB5) are independently replayed in
\path{check_prime_harmonic_cartans.py}; its exact
integer and finite-field outputs are archived.

The literal frame composition law of R1 E.65.12
and the covariance give
$W_{p,hh'}=W_{p,h}W_{p,h'}$, so $w_p$ is a
continuous character.  For an arbitrary $g\in G$,
the full intertwiner of E.65.15 and the same
frame composition law give
\[
 W_{p,ghg^{-1}}
 \equiv \mathcal W_{p,g}W_{p,h}\mathcal W_{p,g}^{-1}
                           \pmod{\{\pm I\}}.
 \tag{B.AC11}
\]
The projective matrix $\mathcal W_{p,g}$ centralizes
the projective source involution.  Any matrix lift
therefore conjugates $S$ to $S$ or $-S$ and
normalizes $T_p$.  Conjugation preserves the unique
index-two subgroup of this cyclic group.
The sign ambiguity in (B.AC11) has zero character
by $\varepsilon_p(-I)=0$.  This proves
$w_p(ghg^{-1})=w_p(h)$.
Since $D$ is an actual homomorphism, the same
conjugation invariance holds for $d_p$.

For $h\in H\cap G_{\mathbb Q}$ the actual arithmetic
comparison gives $Dh=chc$, hence $d_p(h)=0$.
The compact semidirect-product quotient argument of
B.2.26 therefore supplies the unique full character
\[
 \overline d_p(g)=d_p(g\gamma^{-1}),\qquad
       \gamma\in G_{\mathbb Q},\quad
                          \lambda(\gamma)=\lambda(g).
 \tag{B.AC12}
\]
For $h\in H\cap\widehat{GT}$ the same equality
$Dh=chc$ follows from B.1.8 and gives $d_p(h)=0$.
Arithmetic correction in (B.AC12) proves full GT
vanishing.  Also
$w_p(Dh)=w_p(D_Lh)$ by conjugation invariance.
Thus $d_p$ is zero on $C_L\cap H$.
The exact factorization
$C_L=(C_L\cap H)G_{\mathbb Q}$ in B.2.28 proves
full $C_L$ vanishing.
For $p=3$, equation (B.C24) identifies $d_3$ in
(B.AC4) with the earlier $d_2$ of (B.Y1).
Uniqueness of the full arithmetic-vanishing extension
therefore gives $\overline d_3=\overline d_2$.

We construct the independent kernel values.
The stated set of primes is infinite by an elementary
argument.  Given any finite set of odd primes, choose
a positive odd integer $N$ which is three modulo eight
and one modulo each prime in that set, also imposing
one modulo three.  The character of
$(\mathbb Z/8)^\times$ which is negative on three
and five is multiplicative.  Its value on $N$ is
negative, so some prime divisor of $N$ is three
or five modulo eight and occurs to an odd exponent.
That prime is outside the prescribed set and is
at least five.  This proves infinitude and makes the
ordered sequence in the statement unambiguous.

For $i\geq1$ put
\[
 m_i=12\prod_{j<i}p_j,\qquad
 R_i=B/\operatorname{core}_B
           \ker\!\left(\vartheta_{m_i}:F\to
                   \operatorname{SL}_2(\mathbb Z/m_i)\right),
 \tag{B.AC13}
\]
where $\vartheta_{m_i}$ has the same two integral
generator matrices as (B.AC1).
This is an actual finite source quotient retaining
the signed pure evaluation modulo $m_i$.
Induction from the index-six subgroup $F$ embeds
its image into
$\operatorname{SL}_2(\mathbb Z/m_i)^6\rtimes S_3$.
In particular its order divides
$6|\operatorname{SL}_2(\mathbb Z/m_i)|^6$.
Counting the first nonzero column and then the
second column, and lifting through successive
prime powers, gives
\[
 |\operatorname{SL}_2(\mathbb Z/m)|
       =\prod_{\ell^a\parallel m}\ell^{3a-2}(\ell^2-1).
 \tag{B.AC14}
\]
Every prime divisor of this expression for $m_i$
is at most the largest of $3,p_1,\ldots,p_{i-1}$.
Indeed a prime divisor of $\ell-1$ is less than
$\ell$, and a prime divisor of the even number
$\ell+1$, for odd $\ell$, is at most
$(\ell+1)/2\leq\ell$.
Thus $p_i\nmid|R_i|$, as required by B.2.29.

Apply that result to $(R_i,p_i)$ and let $a_i$
be its full normalized element in $\ker D$.
Before normalization its pure frames have
$\vartheta_{m_i}(\beta)=\vartheta_{m_i}(A)=I$,
while their new projective values are
$q_{p_i}(\beta)=[S]$, $q_{p_i}(A)=1$.
The matrix effect of normalization must be retained,
because its exponent $t$ need not be divisible by
$p_i$.  Write $X_p=\vartheta_p(x)$.
The new frames are
$\beta_0=x^t\beta$, $A_0=x^tA$.
At every retained source representation, the
unnormalized action on the signed braid generators
is the identity: its two conjugating matrices are
$I,I$ at the old factors, and $\pm S,\pm I$ at
the new factor, respectively.  Therefore after
normalization its covariance is exactly
$\operatorname{Ad}(X_p^t)$.
The unique normalized upper-unipotent intertwiner
is consequently $\mathcal V_{p,a_i}=X_p^t$.
It follows that
\[
 W_{p,a_i}
  =(X_p^t\vartheta_p(\beta))^{-1}X_p^t
  =\vartheta_p(\beta)^{-1}.
 \tag{B.AC15}
\]
Both factors left-multiply by the same actual
matrix; no exponent is discarded.

At each earlier prime this value is $I$.
At the new prime it is $\pm S^{-1}$, so (B.AC10)
gives $w_{p_i}(a_i)=1$.
The same calculation modulo twelve gives both
original signed frames $T=W=I$: the two raw
signed pure frame values there were exactly $I$,
by (B.AC13).  Hence $\Delta(a_i)=0$.
Since $Da_i=1$, (B.C24) gives
$\overline d_2(a_i)=\epsilon_2(W_{a_i})=0$.
For every new odd-prime character, restriction
to $\ker D$ is simply $d_p=w_p$.
We have therefore proved the exact triangular
kernel values
\[
 \Delta(a_i)=\overline d_2(a_i)=0,\qquad
 \overline d_{p_j}(a_i)=0\ (j<i),\qquad
                         \overline d_{p_i}(a_i)=1.
 \tag{B.AC16}
\]
Together with the two kernel vectors of B.2.27,
these imply independence of all restrictions
$\kappa_j|_{\ker D}$: evaluate any finite relation
at the largest prime-indexed witness occurring,
and then at the two original kernel vectors.

Surjectivity of $D:G\to G$ makes $D^*$ injective.
For any finite polynomial relation in (B.AC6),
evaluation on $\ker D$ makes every constant
coefficient zero by the just-proved independence.
The relation is therefore divisible by $t$;
injectivity of $D^*$ removes that factor.
Repeating proves every coefficient zero.
The same argument works on $H$, since $D:H\to H$
is onto and the kernel witnesses lie in $H$.

There is also the exact image identity
\[
 D^*V
    =\{\eta\in V:\eta|_{\ker D}=0\}.
 \tag{B.AC17}
\]
Indeed a character on the right factors uniquely
and continuously through the quotient map $D$.
Its descended character still vanishes on
$\widehat{GT}$, because $D$ maps that subgroup
onto itself by its specified inner action.
The converse follows directly from $D(\ker D)=1$.
Thus restriction identifies $V/D^*V$ with its
actual image on $\ker D$.  The independent
restrictions (B.AC16) make its dimension at least
$\aleph_0$, and the countability of $V$ in B.2.25
gives the matching upper bound.  This proves
(B.AC7).  A finitely generated $\mathbb F_2[t]$
module has a finite-dimensional quotient by $t$,
so the asserted failure of finite generation follows.

Finally independence in (B.AC6), already on $H$,
makes every finite coordinate projection of
(B.AC8) surjective: a proper image would admit
a nonzero annihilating linear functional.
Compactness of $H$ supplies the inverse image
of every infinite coordinate pattern.  This
proves its full surjectivity and all finite
image assertions.  Its shift relation is
the defining equality
$\kappa_j(D^nDg)=\kappa_j(D^{n+1}g)$.
Each coordinate vanishes on GT, since $D$
preserves that subgroup.  For $g\in C_L$,
equation (B.AA10) gives $D^ng=c^ngc^n$;
the full character vanishing already proved
therefore also holds for every $j\geq1$ and
all $n$.  This completes the proof.
\end{proof}

The free submodule in (B.AC6) has a countably infinite
specified set of generators.  Its identification does
not assert that it exhausts $V$; (B.AC7) is the exact
quotient statement about the entire annihilator.

